\subsection{Enunciado}

\begin{enunciado}{Tarjetas ABC}
La demanda final es de \(200\) tarjetas/hora. El proceso tiene tres operaciones antes del control de calidad: \(P1\), \(P2\) y \(P3\).

\begin{center}
\begin{tabular}{ccccc}
\toprule
Operacion & Lote \(C\) & \(t_p\) s & \(t_k\) s & \(t_v\) s \\
\midrule
P1 & 200 & 85 & 45 & 200 \\
P2 & 250 & 78 & 67 & 300 \\
P3 & 300 & 50 & 92 & 150 \\
\bottomrule
\end{tabular}
\end{center}

Control de calidad descarta \(15\%\). Se pide calcular Kanbans. Luego se estudia P3 con cinco muestras de largos:

\begin{center}
\begin{tabular}{cccccc}
\toprule
Dia & \multicolumn{5}{c}{Largo (mm)} \\
\midrule
1 & 70 & 56 & 49 & 67 & 61 \\
2 & 65 & 47 & 70 & 70 & 68 \\
3 & 70 & 49 & 42 & 68 & 54 \\
4 & 50 & 47 & 52 & 67 & 50 \\
5 & 48 & 65 & 51 & 50 & 65 \\
\bottomrule
\end{tabular}
\end{center}

Se pide calcular limites eliminando outliers si corresponde, evaluar la muestra del dia 6 \((63,54,43,69,65)\), y determinar si cambia el numero de Kanbans cuando las fallas bajan a \(5\%\).
\end{enunciado}
